\begin{table}[t]\centering\small\caption{模型容量与代码来源。AST-BPE采用公开算法的数学迁移实现。}\label{tab:models}\begin{tabular}{lrr}\toprule 模型&ID空间&来源\\\midrule
AST-BPE & 1,024 & \href{https://github.com/SamuelGabriel/sqlbpe}{代码}\\
软AST-BPE & 1,024 & \href{https://github.com/xtr1976536/math-tokenizer-lab}{代码}\\
数学BPE & 1,024 & \href{https://github.com/huggingface/tokenizers}{代码}\\
BPE & 1,024 & \href{https://github.com/huggingface/tokenizers}{代码}\\
Unigram & 1,024 & \href{https://github.com/google/sentencepiece}{代码}\\
Qwen2.5 & 151,665 & \href{https://huggingface.co/Qwen/Qwen2.5-0.5B/blob/main/tokenizer.json}{代码}\\
Llama 3 & 128,256 & \href{https://github.com/meta-llama/llama3/blob/main/llama/tokenizer.py}{代码}\\
o200k & 200,019 & \href{https://github.com/openai/tiktoken}{代码}\\
\bottomrule\end{tabular}\end{table}
